\documentclass[10pt]{article}
\usepackage[T1]{fontenc}
\usepackage[utf8]{inputenc}
\usepackage[french]{babel}
\usepackage{amsmath,amssymb,mathrsfs,mathpazo}
\usepackage{geometry}
\geometry{a4paper,margin=2cm}
\usepackage{graphicx,array,tabularx,tabularray,colortbl,multirow,multicol,booktabs,enumitem,pifont,ulem,fancybox,diagbox,cancel}
\usepackage{pstricks,pstricks-add,pst-plot,pst-tree,pst-node,pst-text,pst-func,pst-math,pst-bspline,pst-3dplot,pst-eucl,pst-solides3d}
\let\pscalculate\undefined
\usepackage{pst-bezier}
\usepackage{tikz,pgfplots,tkz-tab}
\usetikzlibrary{arrows,arrows.meta,patterns,calc,positioning,trees,decorations.markings,shapes.geometric,shapes.misc,shapes.arrows,intersections,decorations.pathmorphing,decorations.pathreplacing,backgrounds,shadows,babel}
\pgfplotsset{compat=1.18}
\usepgfplotslibrary{fillbetween}
\usepackage[np]{numprint}
\usepackage{listings,makecell,pas-tableur,twoopt}
\let\barre\undefined
\usepackage[pstricks]{tablvar}
\definecolor{darkgreen}{RGB}{0,100,0}
\providecommand{\red}{\color{red}}
\providecommand{\blue}{\color{blue}}
\providecommand{\green}{\color{green}}
\providecommand{\black}{\color{black}}
\providecommand{\white}{\color{white}}
\providecommand{\gray}{\color{gray}}
\providecommand{\cyan}{\color{cyan}}
\providecommand{\magenta}{\color{magenta}}
\providecommand{\yellow}{\color{yellow}}
\providecommand{\violet}{\color{violet}}
\providecommand{\orange}{\color{orange}}
\providecommand{\euro}{\texteuro}
\providecommand{\EUR}{\texteuro}
\providecommand{\ieme}{\textsuperscript{e}}
\providecommand{\degres}{\ensuremath{^\circ}}
\providecommand{\degree}{\ensuremath{^\circ}}
\providecommand{\vv}[1]{\overrightarrow{#1}}
\providecommand{\Coord}[1]{#1}
\providecommand{\decosix}{}
\providecommand{\starredbullet}{\ensuremath{\bullet}}
\pagestyle{empty}
\setlength{\parindent}{0pt}
\usepackage{adjustbox}\DeclareUnicodeCharacter{2236}{:}\DeclareUnicodeCharacter{2190}{\ensuremath{\leftarrow}}
\providecommand{\R}{\mathbb{R}}
\providecommand{\N}{\mathbb{N}}
\providecommand{\Z}{\mathbb{Z}}
\providecommand{\Q}{\mathbb{Q}}
\providecommand{\D}{\mathbb{D}}
\providecommand{\C}{\mathbb{C}}
\providecommand{\e}{{\mathrm{e}}}
\providecommand{\vect}[1]{\overrightarrow{#1}}
\providecommand{\vectt}[1]{\overrightarrow{\mathrm{#1}}}
\providecommand{\vvt}[1]{\overrightarrow{\mathrm{#1}}}
\providecommand{\np}[2][]{{#2~#1}}
\providecommand{\vv}[1]{{\overrightarrow{#1}}}
\renewcommand{\d}{{\mathrm{d}}}
\providecommand{\degres}{{^\circ}}
\providecommand{\degre}{{^\circ}}
\providecommand{\blue}{\color{blue}}
\providecommand{\red}{\color{red}}
\providecommand{\enskip}{\quad}
\providecommand{\llbracket}{\lbrack\!\lbrack}
\providecommand{\rrbracket}{\rbrack\!\rbrack}
\providecommand{\euro}{{\text{€}}}
\providecommand{\no}{{\text{n°}}}
\providecommand{\interval}[2]{{[#1;#2]}}
\providecommand{\intervFO}[2]{{[#1;#2[}}
\providecommand{\dotfill}{\ldots}

\providecommand{\ds}{\displaystyle}
\providecommand{\textding}[1]{#1}
\providecommand{\Oij}{\(\mathrm{O};\vec{\imath},\vec{\jmath}\)}
\providecommand{\Oijk}{\(\mathrm{O};\vec{\imath},\vec{\jmath},\vec{k}\)}
\providecommand{\Ouv}{\(\mathrm{O};\vec{u},\vec{v}\)}
\providecommand{\intervFF}[2]{\ensuremath{[#1;#2]}}
\renewcommand{\euro}{\ifmmode\text{\texteuro}\else\texteuro\fi}
\DeclareUnicodeCharacter{00B0}{\ifmmode^{\circ}\else\textdegree\fi}
\DeclareUnicodeCharacter{20AC}{\ifmmode\text{\texteuro}\else\texteuro\fi}
% Énoncés : Éducation nationale. Transcription et corrigés : APMEP, Denis Vergès et contributeurs. Découpage et classement : CoopMaths, Sébastien Lozano, Éric Elter et contributeurs. Adaptation de la présentation : Maths Explorer.
% CC BY-SA : https://coopmaths.fr/www/about/
% https://coopmaths.fr/alea/static/e3c/2020/tex/e3c_2020_00_sujet28_4.tex
\begin{document}
\noindent\textit{APMEP / CoopMaths. CC BY-SA. Présentation adaptée par Maths Explorer.}\par\medskip
\medskip

On procède, chez un sportif, à l'injection intramusculaire d'un produit.

Celui-ci se diffuse progressivement dans le sang. On admet que la concentration de ce produit dans le sang (exprimée en mg/L = milligramme par litre) peut être modélisée par la fonction $f$, définie sur l'intervalle [0~;~10] par :

\[f(x)=\dfrac{6x}{\e^x} \]

où $x$ est le temps exprimé en heure.

Sa courbe représentative  est donnée ci-dessous dans un repère orthonormé du plan.

\psset{xunit=1.2cm,yunit=1.2cm,labelFontSize=\scriptstyle,showorigin=false}
\begin{pspicture}(-1,-1.4)(10.5,3.7)
\multido{\n=-0.2+0.2}{56}{\psline[linewidth=0.35pt,linecolor=lightgray](\n,-1)(\n,3)}
\multido{\n=-1.0+0.2}{21}{\psline[linewidth=0.35pt,linecolor=lightgray](0,\n)(10.8,\n)}
\multido{\n=0+1}{11}{\psline[linewidth=0.45pt](\n,-1)(\n,3)}
\multido{\n=-1+1}{5}{\psline[linewidth=0.45pt](0,\n)(10.8,\n)}
\psaxes[linewidth=0.95pt]{->}(0,0)(0,-1)(11,3)
\def\Func{x 6 mul  2.71828 x exp div}
\psplot[plotpoints=3000,linewidth=1.25pt,linecolor=red]{0}{10}{\Func}
\uput[ur](2.5,1.5){\red $\mathcal{C}$}
\end{pspicture}

\begin{enumerate}
\item Montrer que pour tout $x\in [0~;~10]$, la fonction dérivée de $f$, notée $f'$, a pour expression : 

\[f'(x)=\dfrac{6-6x}{\e^x}.\]

\item Étudier le signe de $f'(x)$ sur [0~;~10] puis en déduire le tableau de variations de $f$ sur [0~;~10].
\item Quelle est la concentration maximale du médicament dans le sang ? 

\emph{(On donnera la valeur exacte et une valeur approchée à $10^{-1}$ près).}

Au bout de combien de temps est-elle atteinte ?
\item Ce produit fait l'objet d'une réglementation par la fédération sportive : un sportif est en infraction si, au moment du contrôle, la concentration dans son sang du produit est supérieure à  \np[mg/L]{2}.

Le sportif peut-il être contrôlé à tout moment après son injection ?

 Expliquer votre raisonnement en vous basant sur l'étude de la fonction et/ou une lecture graphique sur la courbe $\mathcal{C}$.
\end{enumerate}

\end{document}
